\chapter{Cohomological Field Theories}
\label{ch:tqft-cohomological-field-theories}

\ConventionsEuclidean

A cohomological field theory is a quantum field theory equipped with an odd
symmetry \(Q\) whose cohomology controls a protected class of observables.
Its data include the regulated fields, the algebra on which \(Q\) acts, the
closure relation for \(Q^2\), the integration cycle or state, and the
anomaly-free \(Q\)-Ward identity.  Once those entries are supplied, metric
independence and localization are consequences of Stokes' theorem in the
regulated field space or of the BV quantum master equation.

This chapter gives the abstract cohomological mechanism in a form usable by
the later Donaldson--Witten, topological sigma-model, and twisted
supersymmetric examples.  The finite-dimensional statements are proved
directly.  Infinite-dimensional QFT examples are presented as constructions
that must be implemented at regulator level; the later chapters identify
which parts are theorem-level mathematics and which parts still depend on
unproved continuum-QFT input.

\section{Regulated Cohomological Datum}

Let \(\Lambda\) denote a finite regulator: a lattice, a finite-dimensional
mode cutoff, a finite-dimensional approximation to a BV complex, or any other
scheme for which the integral is an actual finite-dimensional integral or a
finite trace.  The regulator is part of the notation because the Ward
identity is false until it has been proved at some regulator and shown to
survive the continuum limit.

\begin{definition}[Regulated cohomological QFT data]
\label{def:regulated-cohomological-qft-data}
Regulated cohomological QFT data consist of:
\begin{enumerate}[label=(\roman*)]
\item a \(\mathbb Z_2\)-graded algebra
  \(\mathcal A_\Lambda\) of regulated insertions;
\item an odd derivation
  \[
    Q_\Lambda:\mathcal A_\Lambda\to\mathcal A_\Lambda;
  \]
\item a distinguished even element \(S_\Lambda\), or more generally a
  regulated state/weight, defining expectation functionals
  \[
    \langle X\rangle_\Lambda
    =
    Z_\Lambda^{-1}\int_{\Gamma_\Lambda} X\,\exp(-S_\Lambda);
  \]
\item a declared integration cycle \(\Gamma_\Lambda\), contour, BV
  Lagrangian, or Hilbert-space trace;
\item an even derivation \(B_\Lambda\) such that
  \begin{equation}
    Q_\Lambda^2=B_\Lambda
    \label{eq:cohomological-q-square-b}
  \end{equation}
  on \(\mathcal A_\Lambda\);
\item a \(B_\Lambda\)-invariant subalgebra
  \[
    \mathcal A_\Lambda^{B}
    =
    \ker B_\Lambda
  \]
  on which \(Q_\Lambda^2=0\);
\item the Ward identity
  \begin{equation}
    \langle Q_\Lambda Y\rangle_\Lambda=0
    \label{eq:regulated-q-ward-identity}
  \end{equation}
  for every admissible insertion \(Y\) used in the argument.
\end{enumerate}
The regulated cohomological observables are the classes
\[
  H_{Q_\Lambda}(\mathcal A_\Lambda^{B})
  =
  \frac{\ker(Q_\Lambda:\mathcal A_\Lambda^B\to\mathcal A_\Lambda^B)}
       {\operatorname{im}(Q_\Lambda:\mathcal A_\Lambda^B\to\mathcal A_\Lambda^B)} .
\]
\end{definition}

The even derivation \(B_\Lambda\) may be zero.  In important examples it is
the infinitesimal action of a compact bosonic symmetry, a gauge
transformation, a flavor transformation, or a combination of these.  The
restriction to \(\mathcal A_\Lambda^B\) is the Cartan-model step: \(Q\) is a
cohomological differential on invariant observables even when it is not
nilpotent on all fields.

If the data are presented by a finite-dimensional supermanifold, the odd
coordinates are coordinates in the structure sheaf of a locally
super-ringed space; they are not points of an ordinary set.  Berezin
integration is the algebraic operation extracting the top odd coefficient,
possibly after multiplication by a Berezinian density.  It is not a Borel
measure.  This distinction is harmless in finite-dimensional algebraic
manipulations but becomes important when one discusses continuum
``fermionic field configurations.''

\begin{definition}[Anomaly and boundary term]
\label{def:cohomological-anomaly-boundary}
If instead of \eqref{eq:regulated-q-ward-identity} one has
\[
  \langle Q_\Lambda Y\rangle_\Lambda
  =
  \langle \mathcal A_\Lambda(Y)\rangle_\Lambda
  +
  \langle \partial_\Lambda(Y)\rangle_\Lambda,
\]
then \(\mathcal A_\Lambda\) is the local cohomological anomaly and
\(\partial_\Lambda\) is the boundary contribution from the integration
cycle, compactification, contour endpoint, or singular stratum.  A
cohomological statement is anomaly-free for the chosen observables only when
both terms vanish or cancel by a stated inflow or boundary prescription.
\end{definition}

\section{Ward Identities and Metric Independence}

Let \(g\) be a background metric, and write
\(\delta_g\) for a tangent vector to the space of metrics.  A metric
variation formula is a statement in a fixed renormalization and contact-term
scheme:
\begin{equation}
  \delta_g\langle X\rangle_\Lambda
  =
  \left\langle
    \delta_g X-(\delta_g S_\Lambda)X
  \right\rangle_\Lambda
  -
  \langle X\rangle_\Lambda\langle-\delta_gS_\Lambda\rangle_\Lambda,
  \label{eq:cohomological-normalized-metric-variation}
\end{equation}
with additional contact insertions included in \(\delta_gX\) when the
representative of \(X\) depends on the metric.

\paragraph{Metric-independence criterion.}
Let \(X_g\in\mathcal A_\Lambda^B\) be a smooth family of
\(Q_\Lambda\)-closed insertions.  Suppose the normalized metric variation can
be written as
\begin{equation}
  \delta_g\langle X_g\rangle_\Lambda
  =
  \langle Q_\Lambda Y_{\delta g}\rangle_\Lambda
  \label{eq:metric-variation-q-exact-insertion}
\end{equation}
for an admissible insertion \(Y_{\delta g}\).  If the Ward identity
\eqref{eq:regulated-q-ward-identity} holds for \(Y_{\delta g}\), then
\[
  \delta_g\langle X_g\rangle_\Lambda=0.
\]
Indeed, equation~\eqref{eq:metric-variation-q-exact-insertion} identifies the
metric derivative with \(\langle Q_\Lambda Y_{\delta g}\rangle_\Lambda\), and
the regulated Ward identity evaluates this expectation value to zero.  The
substantive part of using the criterion is constructing \(Y_{\delta g}\) with
all contact terms and metric dependence of \(X_g\) included in the regulated
insertion algebra.

The familiar condition that the stress tensor is \(Q\)-exact is a sufficient
local mechanism, not the full statement.  In the notation of
Chapter~\ref{ch:tqft-metric-independence}, the complete condition is that
the full renormalized metric response, including contact terms and explicit
variation of the insertion representative, is a \(Q\)-variation.  This is
the condition expressed by
\eqref{eq:metric-variation-q-exact-insertion}.

\section{Exact Deformations and Localization}

Let \(V\in\mathcal A_\Lambda\) be odd and \(B_\Lambda\)-invariant.  If
\begin{equation}
  S_{\Lambda,t}=S_\Lambda+t\,Q_\Lambda V
  \label{eq:q-exact-deformed-action}
\end{equation}
is an allowed family of regulated actions and \(X\) is
\(Q_\Lambda\)-closed, then \(t\)-independence is the same Ward identity.

\paragraph{\(Q\)-exact deformation independence.}
Assume \(Q_\Lambda S_\Lambda=0\), \(Q_\Lambda X=0\),
\(Q_\Lambda^2V=0\), and the Ward identity holds for
\[
  Y=V X\exp(-tQ_\Lambda V)
\]
for every \(t\) in the interval under consideration.  Then
\[
  \frac{\dd}{\dd t}\langle X\rangle_{\Lambda,t}=0 .
\]

For the unnormalized integral
\[
  I_t(X)=\int_{\Gamma_\Lambda}X\exp(-S_\Lambda-tQ_\Lambda V),
\]
one has
\[
  \frac{\dd}{\dd t}I_t(X)
  =
  -\int_{\Gamma_\Lambda}X(Q_\Lambda V)\exp(-S_\Lambda-tQ_\Lambda V).
\]
Since \(Q_\Lambda X=0\), \(Q_\Lambda S_\Lambda=0\), and
\(Q_\Lambda^2V=0\), the integrand is
\[
  -Q_\Lambda\!\left(VX\exp(-S_\Lambda-tQ_\Lambda V)\right).
\]
The Ward identity makes the derivative of \(I_t(X)\) vanish.  Applying the
calculation above with \(X=1\) gives
\[
  \frac{\dd}{\dd t}I_t(1)
  =
  -\int_{\Gamma_\Lambda}
  Q_\Lambda\!\left(V\exp(-S_\Lambda-tQ_\Lambda V)\right)
  =
  0,
\]
so the normalization is also \(t\)-independent.  The quotient
\(\langle X\rangle_{\Lambda,t}=I_t(X)/I_t(1)\) is therefore independent of
\(t\).

This calculation is the algebraic part of localization.  It does not by
itself identify the locus, the determinant, or the boundary contribution.
Those require positivity or contour information for the bosonic part of
\(Q_\Lambda V\), a compactness statement, and a normal-mode analysis.

\section{Finite-Dimensional Localization Theorem}

The following theorem is the precise finite-dimensional model used by
cohomological path integrals.  It is intentionally stated with hypotheses
that are checkable in finite dimension.

\begin{theorem}[Compact nondegenerate localization]
\label{thm:compact-nondegenerate-cohomological-localization}
Let \(\mathcal X\) be an oriented finite-dimensional supermanifold with
compact even integration cycle \(\Gamma\), Berezinian density \(\mu\), and an
odd vector field \(Q\).  Let \(S\in C^\infty(\mathcal X)_{\bar 0}\) be an
even action, \(X\in C^\infty(\mathcal X)\) an insertion, and
\(V\in C^\infty(\mathcal X)_{\bar 1}\) an odd deformation functional.
Suppose:
\begin{enumerate}[label=(\roman*)]
\item \(Q^2\) generates a compact even symmetry whose invariant subalgebra is
  used for observables;
\item \(L_Q\mu=0\) and \(Q\) has no boundary flux through \(\partial\Gamma\);
\item \(QS=0\), \(QX=0\), and \(S\) and \(X\) are invariant under \(Q^2\);
\item \(V\) is invariant under \(Q^2\);
\item the even part \(h=(QV)_{\rm bos}\) satisfies \(h\ge0\), and
  \(F=h^{-1}(0)\subset\Gamma_{\rm red}\) is a compact smooth submanifold;
\item in a tubular neighborhood of \(F\), the normal expansion of \(QV\)
  has a nondegenerate quadratic normal super-Gaussian whose local inverse
  Euler factors patch to a smooth density-valued section
  \(e_{Q^2}(N_F)^{-1}\) on \(F\).
\end{enumerate}
More explicitly, the last condition means the following.  For each
\(y\in F\), after choosing even normal coordinates \(u\in N_{\bar 0,y}\) and
odd normal coordinates \(\eta\in N_{\bar 1,y}\), the quadratic normal part is
\begin{equation}
  (QV)_2(y;u,\eta)
  =
  \frac12 u^iA_{ij}(y)u^j
  +
  \frac12 \eta^\alpha C_{\alpha\beta}(y)\eta^\beta ,
  \label{eq:cohomological-normal-quadratic-model}
\end{equation}
where \(A_y\) is positive definite on the chosen real normal contour and
\(C_y\) is nondegenerate on the odd normal variables.  The leading
Berezinian degree is assumed to be balanced: after the normal rescaling
\(u=t^{-1/2}\xi\), \(\eta=t^{-1/2}\zeta\), the product of the even Jacobian,
the odd Berezinian, and the degree of the normal Berezin polynomial has a
finite \(t^0\) limit.  If this balance fails, the right large-\(t\) statement
is a different asymptotic expansion rather than the simple fixed-locus
formula below.  Under this balance condition, the normal Gaussian
\begin{equation}
  \mathcal G(y)
  =
  \int_{N_y}
  \exp\!\left[-(QV)_2(y;u,\eta)\right]\,D u\,D\eta
  =
  (2\pi)^{\dim N_{\bar0}/2}
  \frac{\operatorname{Pf}(C_y)}{\det(A_y)^{1/2}}
  \label{eq:cohomological-normal-gaussian-euler-factor}
\end{equation}
with the chosen orientation and square-root convention is the local value of
\(e_{Q^2}(N_F)^{-1}\).  Changes of normal coordinates act on the numerator,
denominator, and Berezinian in the compensating way, so the patched object is
a density on \(F\).  Then
\[
  I_t(X)=\int_{\Gamma}X\exp(-S-tQV)\,\mu
\]
is independent of \(t\ge0\), and its large-\(t\) value is
\begin{equation}
  \int_F
  \frac{X|_F\,\exp(-S|_F)}
       {e_{Q^2}(N_F)}
  \label{eq:cohomological-localization-fixed-locus-formula}
\end{equation}
with the orientation line and equivariant Euler class determined by the
normal Gaussian datum \eqref{eq:cohomological-normal-quadratic-model}.
\end{theorem}

\begin{proof}
The \(t\)-independence is the \(Q\)-exact deformation calculation above in
finite dimension, with \(\langle QY\rangle=0\) following from Stokes'
theorem for the divergence-free odd vector field \(Q\).  In detail,
\[
  \frac{\dd I_t(X)}{\dd t}
  =
  -\int_\Gamma X(QV)e^{-S-tQV}\,\mu
  =
  -\int_\Gamma Q\!\left(VXe^{-S-tQV}\right)\mu
  =
  0,
\]
because \(QS=0\), \(QX=0\), \(Q^2V=0\) on the invariant subalgebra, and the
boundary flux of \(Q\) through \(\partial\Gamma\) vanishes.

It remains to identify the common value by the limit \(t\to+\infty\).  Since
\(\Gamma_{\rm red}\) is compact and \(h^{-1}(0)=F\), for every sufficiently
small tubular neighborhood \(U\) of \(F\) there is a constant \(c_U>0\) with
\(h\ge c_U\) on \(\Gamma_{\rm red}\setminus U\).  The even body of
the integrand outside \(U\) is therefore bounded by
\(\exp(-tc_U)\) times a fixed smooth density, and the odd part is a finite
polynomial.  Hence the complement of \(U\) does not contribute to the limit.

Work in one coordinate patch of the tubular neighborhood and write a point as
\((y,u,\eta)\), with \(y\in F\), even normal coordinate \(u\), and odd normal
coordinate \(\eta\).  Taylor expansion gives
\[
  QV=(QV)_2+R_3,
\]
where \((QV)_2\) is the quadratic expression
\eqref{eq:cohomological-normal-quadratic-model} and \(R_3\) has total normal
degree at least three.  Put \(u=t^{-1/2}\xi\) and
\(\eta=t^{-1/2}\zeta\).  In Berezin integration the odd change of variables
contributes the inverse determinant of the odd scaling, while the even
change of variables contributes the ordinary Jacobian; these are exactly the
factors included in the super-Gaussian density
\eqref{eq:cohomological-normal-gaussian-euler-factor}.  The remainder
satisfies
\[
  t\,R_3(y,t^{-1/2}\xi,t^{-1/2}\zeta)=O(t^{-1/2})
\]
on every bounded normal set, and the positive quadratic form \(A_y\) gives a
uniform Gaussian majorant after shrinking the tubular neighborhood and using
a finite cover of \(F\).  Dominated convergence therefore replaces the
normal integral by the fiberwise Gaussian \(\mathcal G(y)\).

The local Gaussian factors transform by the Berezinian of changes of normal
coordinates, so they glue to the inverse Euler density
\(e_{Q^2}(N_F)^{-1}\).  Summing the finite partition of unity on \(F\)
gives \eqref{eq:cohomological-localization-fixed-locus-formula}.  The
orientation line and the square root of \(\det A_y\) are not decoration:
without these choices the local Gaussian factors do not define a global
integrand on \(F\).
\end{proof}

\section{Mathai--Quillen Form of a Cohomological Theory}

Many cohomological field theories are infinite-dimensional analogues of the
Mathai--Quillen representative of the Euler class of a vector bundle.  The
finite-dimensional model is as follows.

Let \(\pi:E\to X\) be an oriented real vector bundle with metric
\((\cdot,\cdot)_E\), compatible connection \(\nabla^E\), and smooth section
\[
  s:X\to E.
\]
The zero locus is
\[
  Z(s)=\{x\in X:s(x)=0\}.
\]
Introduce local even coordinates \(x^i\), odd tangent variables \(\psi^i\)
representing \(dx^i\), odd fiber variables \(\chi_a\), and even auxiliary
fiber variables \(H_a\).  The cohomological differential is
\begin{equation}
  Qx^i=\psi^i,\qquad Q\psi^i=0,\qquad
  Q\chi_a=H_a,\qquad QH_a=0
  \label{eq:mq-q-rules-flat}
\end{equation}
in a local trivialization.  The connection-covariant version adds the
standard connection terms so that the expression below globalizes.

Choose the standard oscillatory auxiliary-field representative
\[
  V_{\rm MQ}
  =
  \frac12(\chi,H)_E+\ii(\chi,s)_E .
\]
Then
\begin{equation}
  QV_{\rm MQ}
  =
  \frac12(H,H)_E+\ii(H,s)_E
  -\ii(\chi,\nabla_\psi^E s)_E
  +\text{curvature term}.
  \label{eq:mq-qv}
\end{equation}
With the real Gaussian contour for \(H\), completing the square gives
\[
  \int \dd H\,
  \exp\left[-\frac12(H,H)_E-\ii(H,s)_E\right]
  =
  \text{const.}\,
  \exp\!\left[-\frac12\|s\|^2\right].
\]
The remaining Berezin polynomial, including the connection and curvature
terms suppressed in \eqref{eq:mq-qv}, has top fiber-degree component equal to the
Mathai--Quillen Thom form.  Pulling this Thom form back by the section
\(s\) gives a de Rham representative of \(e(E)\).  If the section is
transverse, the large-deformation limit localizes the integral to \(Z(s)\),
with sign determined by the orientation of \(\det(ds)\).

\begin{proposition}[Mathai--Quillen localization for a transverse section]
\label{prop:mq-zero-locus-interpretation}
Let \(X\) be compact and \(s\) transverse to the zero section.  For any
closed differential form \(\alpha\) on \(X\) of degree
\(\dim X-\operatorname{rank}E\),
\[
  \int_X \alpha\wedge e(E)
  =
  \int_{Z(s)} \alpha|_{Z(s)}
  \]
where \(Z(s)\) is oriented by the isomorphism
\[
  N_{Z(s)/X}\simeq s^*E
\]
induced by \(ds\) and by the given orientations of \(X\) and \(E\).  When
\(\operatorname{rank}E=\dim X\) it is the signed sum of zeros of \(s\).
\end{proposition}

\begin{proof}
Let \(\tau_{\rm MQ}\in\Omega^{\operatorname{rank}E}_{\rm cpt}(E)\) denote the
Mathai--Quillen Thom form in the convention described above.  Its pullback
\(s^*\tau_{\rm MQ}\) represents \(e(E)\).  Replacing \(s\) by \(t^{1/2}s\)
does not change this cohomology class: differentiating the Gaussian
representative with respect to \(t\) gives an exact form, equivalently a
\(Q\)-exact variation in the finite-dimensional supersymmetric notation.
Since \(X\) is compact and \(\alpha\) is closed,
\[
  I(t):=\int_X\alpha\wedge (t^{1/2}s)^*\tau_{\rm MQ}
\]
is independent of \(t>0\).

It remains to identify the common value by the limit \(t\to+\infty\).
Transversality makes \(Z(s)\) a compact submanifold of codimension
\(\operatorname{rank}E\), and
\[
  ds:N_{Z(s)/X}\longrightarrow s^*E
\]
is a vector-bundle isomorphism.  Choose a tubular neighborhood of \(Z(s)\)
with base coordinate \(y\in Z(s)\) and normal coordinate \(n\in N_{Z(s)/X,y}\).
In this neighborhood
\[
  s(y,n)=ds_y(n)+O(|n|^2).
\]
On the complement of any smaller tubular neighborhood, \(\|s\|\) is bounded
below by a positive constant, so the Gaussian part of the pulled-back
Mathai--Quillen form is exponentially small in \(t\).  Thus only
\(|n|=O(t^{-1/2})\) contributes.  Put \(n=t^{-1/2}\xi\).  The leading normal
part of the Mathai--Quillen density is
\[
  (2\pi)^{-r/2}
  \exp\!\left[-\frac12\|ds_y(\xi)\|^2\right]
  \det(ds_y)\,d^r\xi,
  \qquad r=\operatorname{rank}E,
\]
where the determinant is computed after oriented bases of \(N_{Z(s)/X,y}\)
and \(E_y\) are chosen.  The ordinary Gaussian integral contributes
\(|\det(ds_y)|^{-1}\), while the Berezin fiber integration contributes
\(\det(ds_y)\).  Their quotient is the sign of \(ds_y\).  The orientation on
\(Z(s)\) is exactly the orientation for which the decomposition
\[
  T_xX\simeq T_xZ(s)\oplus N_{Z(s)/X,x}
  \simeq T_xZ(s)\oplus E_x
\]
has the prescribed orientations of \(X\) and \(E\).  With this convention the
sign is absorbed into the oriented integration over \(Z(s)\).  The
terms omitted in \(s(y,t^{-1/2}\xi)\) and in the pulled-back form give
negative powers of \(t^{1/2}\) after the rescaling and vanish by the Gaussian
bound.  Therefore
\[
  \lim_{t\to+\infty} I(t)=\int_{Z(s)}\alpha|_{Z(s)} .
\]
The equality follows because \(I(t)\) is independent of \(t\).  When
\(\operatorname{rank}E=\dim X\), \(Z(s)\) is a finite oriented zero-manifold,
so the last integral is the signed sum over its points.
\end{proof}

The proposition is the finite-dimensional ancestor of statements such as
``Donaldson--Witten theory localizes to anti-self-dual connections'' and
``the A-model localizes to holomorphic maps.''  In field theory, \(X\) is an
infinite-dimensional field space, \(E\) is an infinite-dimensional bundle of
equations, and \(s=0\) is the PDE defining the moduli problem.  The theorem
then requires regularization, compactness, orientation, and determinant-line
control.

\section{Descent and Integrated Observables}

Let \(M\) be a smooth \(D\)-manifold.  A local cohomological observable is
not only a \(Q\)-closed local density.  If it is to be integrated over cycles
of different dimensions, it must come with a descent package.

\begin{definition}[Descent package]
\label{def:descent-package}
A descent package is a sequence of local \(p\)-form insertions
\[
  \mathcal O^{(p)}\in\Omega^p(M)\otimes\mathcal A,
  \qquad
  0\le p\le D,
\]
such that
\begin{equation}
  Q\mathcal O^{(0)}=0,
  \qquad
  Q\mathcal O^{(p)}+\dd\mathcal O^{(p-1)}=0
  \quad (1\le p\le D).
  \label{eq:descent-equations}
\end{equation}
If \(\gamma_p\) is a closed \(p\)-cycle, then
\begin{equation}
  I_{\gamma_p}(\mathcal O)=\int_{\gamma_p}\mathcal O^{(p)}
  \label{eq:descent-integrated-observable}
\end{equation}
is \(Q\)-closed: using \eqref{eq:descent-equations} and
\(\partial\gamma_p=0\),
\[
  Q I_{\gamma_p}(\mathcal O)
  =
  \int_{\gamma_p}Q\mathcal O^{(p)}
  =
  -\int_{\gamma_p}\dd\mathcal O^{(p-1)}
  =
  -\int_{\partial\gamma_p}\mathcal O^{(p-1)}
  =
  0 .
\]
\end{definition}

If \(\gamma_p\) is changed by a boundary,
\(\gamma_p\mapsto\gamma_p+\partial C_{p+1}\), then
\begin{equation}
  I_{\gamma_p+\partial C_{p+1}}(\mathcal O)-I_{\gamma_p}(\mathcal O)
  =
  \int_{C_{p+1}}\dd\mathcal O^{(p)}
  =
  -Q\int_{C_{p+1}}\mathcal O^{(p+1)} .
  \label{eq:descent-homology-q-exact-shift}
\end{equation}
Thus integrated descendants depend only on the homology class of the cycle in
\(Q\)-cohomology, subject to the specified contact-term qualification when
cycles collide with other insertions.

In Donaldson--Witten theory, the descent package is generated from invariant
polynomials in the universal equivariant curvature.  With
\begin{equation}
  QA=\psi,\qquad Q\psi=-d_A\phi,\qquad Q\phi=0,
  \qquad Q^2=\delta_{\rm gauge}(-\phi),
  \label{eq:dw-universal-q-rules}
\end{equation}
gauge-invariant polynomials of \(\phi\) define local \(Q\)-classes, and their
descendants give observables integrated over cycles in the four-manifold.
The detailed four-dimensional field content, instanton moduli space, and
orientation data are developed in
Chapter~\ref{ch:tqft-witten-donaldson-seiberg-witten-comparison}.

\section{Cartan Model and Equivariant Closure}

When \(Q^2\) is an even symmetry, the finite-dimensional model is the Cartan
model of equivariant cohomology.  Let a compact Lie group \(K\) act on a
manifold \(X\) with infinitesimal vector fields \(v_\xi\),
\(\xi\in\mathfrak k\).  On \(K\)-invariant polynomial maps
\[
  \alpha:\mathfrak k\to\Omega^\bullet(X)
\]
define
\begin{equation}
  d_K\alpha(\xi)
  =
  d\alpha(\xi)-\iota_{v_\xi}\alpha(\xi).
  \label{eq:cartan-equivariant-differential}
\end{equation}
Then
\begin{equation}
  d_K^2\alpha(\xi)=-\mathcal L_{v_\xi}\alpha(\xi),
  \label{eq:cartan-square}
\end{equation}
so \(d_K^2=0\) on invariant forms.

\begin{example}[Hamiltonian \(S^1\) model]
\label{ex:hamiltonian-s1-equivariant-closure}
Let \(X=\mathbb R^2\) with coordinates \(x,y\), symplectic form
\(\omega=dx\wedge dy\), and rotation vector field
\[
  v=-y\partial_x+x\partial_y.
\]
The moment map in the sign convention of \eqref{eq:cartan-equivariant-differential}
is
\[
  \mu=-\frac12(x^2+y^2),
  \qquad
  d\mu=\iota_v\omega .
\]
Therefore
\[
  d_{S^1}(\omega+\mu)=0.
\]
This is the elementary algebra behind equivariant localization: the
deformation can be made by an equivariantly exact term without changing the
equivariant cohomology class.
\end{example}

\section{Gauge Theories and BV Closure}

For a gauge theory, replacing \(Q\) by a formal odd symmetry on fields is not
enough.  Gauge fixing, ghosts, antifields, and the integration cycle must be
part of the construction.  The correct compatibility condition is the BV
master equation.  In finite-dimensional notation, let
\((\mathcal F,\omega_{\rm BV})\) be an odd symplectic supermanifold, let
\((\,\cdot\,,\,\cdot\,)\) be the BV bracket, and let \(\Delta\) be the
regulated BV Laplacian.  A quantum BV action satisfies
\begin{equation}
  \frac12(S,S)-\hbar\Delta S=0.
  \label{eq:cohomological-qme}
\end{equation}
The corresponding BV differential on insertions is
\begin{equation}
  \widehat Q X=(S,X)-\hbar\Delta X.
  \label{eq:bv-quantum-differential}
\end{equation}
If \(X=\widehat QY\) and the BV Stokes theorem holds on the chosen
Lagrangian integration cycle \(\mathcal L\), then
\[
  \int_{\mathcal L}X\,\exp(-S/\hbar)=0.
\]

A cohomological gauge theory is therefore a BV theory together with an extra
odd symmetry, or with a BV differential whose cohomology contains the desired
topological observables.  The equality \(Q^2=\delta_{\rm gauge}(\phi)\) must
be interpreted inside the BV complex.  A \(Q\)-exact deformation is
legitimate only if it preserves the quantum master equation, or equivalently
if it is generated by a BV-canonical deformation plus allowed gauge-fixing
data.

\section{Donaldson--Witten as the Guiding Gauge-Theory Model}

The universal finite-dimensional template for Donaldson--Witten theory is
the Mathai--Quillen model applied to the bundle of self-dual two-forms over
the space of connections modulo gauge.  Let \(X\) be a closed oriented
Riemannian four-manifold and \(P\to X\) a principal \(SU(2)\) bundle.  The
configuration space before quotient is the affine space \(\mathcal A(P)\).
The equation bundle is
\[
  \Omega^{2,+}(X,\operatorname{ad}P)\longrightarrow \mathcal A(P),
\]
and the section is
\[
  s(A)=F_A^+ .
\]
The zero locus \(s^{-1}(0)\) is the space of anti-self-dual connections.
After quotient by the gauge group, the virtual tangent-obstruction complex is
\[
  \Omega^0(X,\operatorname{ad}P)
  \xrightarrow{d_A}
  \Omega^1(X,\operatorname{ad}P)
  \xrightarrow{d_A^+}
  \Omega^{2,+}(X,\operatorname{ad}P).
\]
The twisted fields
\[
  \psi\in\Omega^1(X,\operatorname{ad}P),
  \qquad
  \chi^+\in\Omega^{2,+}(X,\operatorname{ad}P),
  \qquad
  H^+\in\Omega^{2,+}(X,\operatorname{ad}P)
\]
are the tangent odd variable, antighost, and auxiliary variable in this
Mathai--Quillen construction.  The scalar \(\phi\) is the equivariant gauge
parameter appearing in \(Q^2=\delta_{\rm gauge}(-\phi)\).  Thus localization
to \(F_A^+=0\) is the infinite-dimensional, gauge-equivariant
Mathai--Quillen localization mechanism, with all analytic and compactness
issues still to be supplied.

This perspective also explains why contact terms are unavoidable.  The
Donaldson observables are descendants integrated over cycles.  Products of
integrated descendants require a prescription for collisions of cycles and
for boundary strata of the instanton compactification.  A cohomological Ward
identity controls these contact terms only after they have been included in
the regulated observable algebra.

\section{What Must Be Proved in Infinite Dimension}

The finite-dimensional statements above reduce the infinite-dimensional
problem to a clear list of mathematical obligations:
\begin{enumerate}[label=(\roman*)]
\item construct the regulated field/BV space and integration cycle;
\item prove the \(Q\)-Ward identity or quantum master equation at the
  regulator;
\item prove that the continuum limit exists for the selected observables;
\item identify the fixed locus and its virtual normal complex;
\item prove compactness or describe boundary strata and their contributions;
\item orient the determinant line and compute the one-loop or Euler factor;
\item define contact terms for integrated descendants and show that the
  descent equations survive renormalization.
\end{enumerate}
Only after these steps are complete does a formal localization computation
become a theorem about a QFT.

\begin{openproblem}[Cohomological QFT beyond formal localization]
\label{op:cohomological-qft-beyond-formal-localization}
Develop regulator-level cohomological QFT constructions in which \(Q\)-Stokes
theorem, BV gauge fixing, compactness of fixed loci, orientation lines,
one-loop determinants, contact terms, and boundary strata are all controlled.
This is the prerequisite for using cohomological field theory as a
theorem-level bridge between quantum field theory and topological invariants.
\end{openproblem}
